\section*{Chapitre \ref{ch:g6:circuits-and-safety} --- Circuits électriques et sécurité}

\begin{solution}{exo:g6:circuits-and-safety:1}
Fidèlement~: les liaisons --- quel appareil rejoint quel autre, où
les routes se séparent. Ignorés~: la longueur, la couleur et les
serpentements réels des fils sur la table.
\end{solution}

\begin{solution}{exo:g6:circuits-and-safety:2}
La \omterm{def:g3:first-electric-circuit:battery}{pile}~: deux barres parallèles inégales~; la lampe~: un cercle
barré d'une croix~; l'\omterm{ex:g3:first-electric-circuit:switch}{interrupteur} ouvert~: un pont-levis relevé~;
l'\omterm{ex:g3:first-electric-circuit:switch}{interrupteur} fermé~: le pont posé à plat~; le \omterm{ex:g6:circuits-and-safety:motor}{moteur}~: un M dans un
cercle~; le \omterm{ex:g6:circuits-and-safety:motor}{vibreur}~: un cercle barré. C'est la barre \emph{longue}
qui marque le $+$ de la \omterm{def:g3:first-electric-circuit:battery}{pile}.
\end{solution}

\begin{solution}{exo:g6:circuits-and-safety:3}
Un seul rectangle~: \omterm{def:g3:first-electric-circuit:battery}{pile}, bouton-poussoir et \omterm{ex:g6:circuits-and-safety:motor}{vibreur} sur une unique
boucle. Retirer le doigt ouvre l'\omterm{ex:g3:first-electric-circuit:switch}{interrupteur} --- une coupure sur
l'unique route --- et la loi de la boucle fait taire le \omterm{ex:g6:circuits-and-safety:motor}{vibreur}
instantanément.
\end{solution}

\begin{solution}{exo:g6:circuits-and-safety:4}
Pas nécessairement. La question qui tranche~: qui est relié à qui~?
Les mêmes appareils avec les mêmes liaisons --- le même circuit,
aussi différents que paraissent les enchevêtrements.
\end{solution}

\begin{solution}{exo:g6:circuits-and-safety:5}
Le ventilateur continue de tourner~: la mort de la lampe n'éteint que
sa propre \omterm{def:g5:series-parallel-circuits:parallel}{branche} --- l'indépendance de la dérivation.
L'\omterm{ex:g3:first-electric-circuit:switch}{interrupteur} ouvert se trouve sur la route commune, avant
l'embranchement~: il les arrête donc tous les deux --- la règle des
postes de commande.
\end{solution}

\begin{solution}{exo:g6:circuits-and-safety:6}
Une route conductrice nue qui relie directement les deux côtés de la
source, en sautant tous les appareils~: rien ne gêne plus le courant,
il se rue et les fils chauffent. Victimes~: la \omterm{def:g3:first-electric-circuit:battery}{pile} (vidée et
fichue), l'\omterm{def:g4:conductors-insulators:insulator}{isolant} du fil --- et, sur le secteur, parfois la maison
entière.
\end{solution}

\begin{solution}{exo:g6:circuits-and-safety:7}
Il a coupé sa branche en un clin d'œil parce que le courant s'est mis
à se ruer --- un défaut, non un caprice. Il faut trouver et réparer
le défaut (débrancher le suspect, inspecter la \omterm{def:g5:series-parallel-circuits:parallel}{branche}) avant de
réarmer le garde.
\end{solution}

\begin{solution}{exo:g6:circuits-and-safety:8}
La baignoire~: la peau humide conduit, et la poussée du secteur à
travers un corps mouillé peut être mortelle. Les cordons abîmés~: un
\omterm{def:g4:conductors-insulators:insulator}{isolant} fendu est une clôture avec un trou --- des doigts peuvent
rencontrer le cuivre. La \omterm{def:g3:first-electric-circuit:battery}{pile} oui, la prise jamais~: la poussée de la
\omterm{def:g3:first-electric-circuit:battery}{pile} est minuscule~; celle de la prise est des centaines de fois plus
puissante.
\end{solution}

\begin{solution}{exo:g6:circuits-and-safety:9}
Schéma~: \omterm{def:g3:first-electric-circuit:battery}{pile}, \omterm{ex:g3:first-electric-circuit:switch}{interrupteur} fermé sur la route commune, puis un
embranchement --- \omterm{def:g5:series-parallel-circuits:parallel}{branche} du \omterm{ex:g6:circuits-and-safety:motor}{vibreur} et \omterm{def:g5:series-parallel-circuits:parallel}{branche} de la lampe --- qui
se rejoignent vers le $-$. Prédiction~: la lampe brille \emph{et} le
\omterm{ex:g6:circuits-and-safety:motor}{vibreur} sonne, indépendamment. Avec un \omterm{ex:g3:first-electric-circuit:switch}{interrupteur} supplémentaire
ouvert sur la \omterm{def:g5:series-parallel-circuits:parallel}{branche} du \omterm{ex:g6:circuits-and-safety:motor}{vibreur}~: le silence de ce côté, mais la
lampe continue de briller --- les \omterm{ex:g3:first-electric-circuit:switch}{interrupteurs} de \omterm{def:g5:series-parallel-circuits:parallel}{branche} ne
commandent que leur \omterm{def:g5:series-parallel-circuits:parallel}{branche}.
\end{solution}

\begin{solution}{exo:g6:circuits-and-safety:10}
Le fil supplémentaire est un \omterm{def:g6:circuits-and-safety:short}{court-circuit}~: une route nue et directe
du $+$ au $-$. Le courant préfère la route vide et facile au \omterm{ex:g6:circuits-and-safety:motor}{moteur}
qui travaille dur~: le ventilateur est donc affamé pendant que le fil
de raccourci encaisse la ruée et chauffe --- ce qui vide la \omterm{def:g3:first-electric-circuit:battery}{pile} et
invite les brûlures. Il faut supprimer le fil «~de solidité~».
\end{solution}

\begin{solution}{exo:g6:circuits-and-safety:11}
L'\omterm{ex:g3:first-electric-circuit:switch}{interrupteur} mural ouvre la boucle propre au luminaire --- la route
de la lampe est coupée. Le disjoncteur ouvre toute la \omterm{def:g5:series-parallel-circuits:parallel}{branche}, bien
plus en amont --- même si le câblage du luminaire est défectueux ou
si l'\omterm{ex:g3:first-electric-circuit:switch}{interrupteur} mural est mal repéré, la \omterm{def:g5:series-parallel-circuits:parallel}{branche} est morte. Deux
coupures \omterm{def:g4:series-circuits:series}{en série}~: chacune arrête le courant à elle seule, et les
deux ensemble pardonnent une erreur sur l'une comme sur l'autre.
\end{solution}

\begin{solution}{exo:g6:circuits-and-safety:12}
La \omterm{def:g3:first-electric-circuit:battery}{pile}, puis l'\omterm{ex:g3:first-electric-circuit:switch}{interrupteur} général sur la route commune, puis
quatre \omterm{def:g5:series-parallel-circuits:parallel}{branches} \omterm{def:g5:series-parallel-circuits:parallel}{en dérivation}~: lampe de scène 1~; lampe de scène 2~;
\omterm{ex:g6:circuits-and-safety:motor}{moteur} de rideau avec son propre \omterm{ex:g3:first-electric-circuit:switch}{interrupteur} \omterm{def:g4:series-circuits:series}{en série}~; \omterm{ex:g6:circuits-and-safety:motor}{vibreur} avec
son propre \omterm{ex:g3:first-electric-circuit:switch}{interrupteur} \omterm{def:g4:series-circuits:series}{en série}. Contrôles~: le général commande
tout (route commune)~; les lampes tombent en panne indépendamment
(\omterm{def:g5:series-parallel-circuits:parallel}{branches} séparées)~; le \omterm{ex:g6:circuits-and-safety:motor}{moteur} et le \omterm{ex:g6:circuits-and-safety:motor}{vibreur} obéissent chacun à son
\omterm{ex:g3:first-electric-circuit:switch}{interrupteur} (\omterm{def:g4:series-circuits:series}{en série} au sein de leur \omterm{def:g5:series-parallel-circuits:parallel}{branche}) et s'ignorent l'un
l'autre (\omterm{def:g5:series-parallel-circuits:parallel}{en dérivation} entre \omterm{def:g5:series-parallel-circuits:parallel}{branches}).
\end{solution}

\begin{solution}{pb:g6:circuits-and-safety:1}
\textbf{1.} \omterm{def:g4:series-circuits:series}{En série}, un seul appareil grillé ou éteint plongerait
tout dans le noir, et les trois appareils se partageraient la
poussée, tous faibles et poussifs. La dérivation garde chaque
appareil indépendant et à pleine puissance.
\textbf{2.} Chaque \omterm{ex:g3:first-electric-circuit:switch}{interrupteur} commande exactement l'appareil de sa
propre \omterm{def:g5:series-parallel-circuits:parallel}{branche}~; aucun d'eux ne commande les autres \omterm{def:g5:series-parallel-circuits:parallel}{branches}.
\textbf{3.} Sur la route commune, entre la \omterm{def:g3:first-electric-circuit:battery}{pile} et le premier
embranchement --- en amont de toutes les \omterm{def:g5:series-parallel-circuits:parallel}{branches}.
\textbf{4.} \omterm{def:g3:first-electric-circuit:battery}{Pile} $\to$ \omterm{ex:g3:first-electric-circuit:switch}{interrupteur} général $\to$ \omterm{def:g5:series-parallel-circuits:parallel}{nœud} d'où partent
trois \omterm{def:g5:series-parallel-circuits:parallel}{branches} --- (plafonnier $+$ son \omterm{ex:g3:first-electric-circuit:switch}{interrupteur}), (lampe
d'établi $+$ son \omterm{ex:g3:first-electric-circuit:switch}{interrupteur}), (\omterm{ex:g6:circuits-and-safety:motor}{moteur} du ventilateur $+$ son
\omterm{ex:g3:first-electric-circuit:switch}{interrupteur}) --- qui se rejoignent et reviennent à la \omterm{def:g3:first-electric-circuit:battery}{pile}.
\textbf{5.} Ils sont \omterm{def:g4:series-circuits:series}{en série} sur une même \omterm{def:g5:series-parallel-circuits:parallel}{branche}~: ils se
partagent la poussée (tous deux poussifs --- le plafonnier terne
était en réalité le symptôme du couple de l'établi), et
l'\omterm{ex:g3:first-electric-circuit:switch}{interrupteur} du ventilateur, placé sur leur route commune, les
commande tous les deux.
\textbf{6.} Les séparer~: tirer un fil neuf pour que la lampe
d'établi ait sa propre \omterm{def:g5:series-parallel-circuits:parallel}{branche} en travers de l'alimentation ---
chacune est alors à pleine puissance et indépendante.
\textbf{7.} Un \omterm{def:g6:circuits-and-safety:short}{court-circuit} entre les \omterm{def:g3:first-electric-circuit:battery}{bornes}. Vérifier si la \omterm{def:g3:first-electric-circuit:battery}{pile}
délivre encore quelque chose --- les courts-circuits la vident et
peuvent la détruire (une \omterm{def:g3:first-electric-circuit:battery}{pile} devenue chaude mérite la retraite).
\textbf{8.} Les courts-circuits (des fils nus qui se rencontrent) et
les chocs ou brûlures (des doigts qui rencontrent le cuivre nu).
\textbf{9.} La peau humide conduit bien mieux que la peau sèche~;
même un montage à \omterm{def:g3:first-electric-circuit:battery}{pile} mérite la discipline des mains sèches --- et,
dans la maison, le même geste rencontre la poussée du secteur, des
centaines de fois plus puissante~: possiblement mortelle. Les
habitudes doivent s'acquérir sur le circuit sans danger.
\textbf{10.} Le disjoncteur a détecté qu'un courant trop grand se
ruait dans sa branche --- un défaut, sans doute de la pluie dans un
appareillage. Ordre raisonnable~: laisser le disjoncteur coupé,
débrancher et inspecter la \omterm{def:g5:series-parallel-circuits:parallel}{branche}, réparer ou sécher le défaut, et
seulement ensuite réarmer.
\textbf{11.} Sa propre boucle, distincte du réseau de la cabane (ou
comme une \omterm{def:g5:series-parallel-circuits:parallel}{branche} \omterm{def:g5:series-parallel-circuits:parallel}{en dérivation} de plus sur la \omterm{def:g3:first-electric-circuit:battery}{pile} de la cabane)~:
\omterm{def:g3:first-electric-circuit:battery}{pile}, bouton-poussoir à la maison, long fil double, \omterm{ex:g6:circuits-and-safety:motor}{vibreur} à la
cabane, et retour --- une sonnette est une boucle \omterm{def:g4:series-circuits:series}{en série} à elle
seule, quelle que soit la longueur de ses fils.
\textbf{12.} Par exemple~: «~Prises~: nos expériences s'arrêtent à la
\omterm{def:g3:first-electric-circuit:battery}{pile} --- la prise n'est pas à nous. Eau~: aucune main mouillée ne
touche un câblage, dedans ou dehors. Fils nus~: chaque jonction
emballée~; un fil nu est un défaut, pas un raccourci.~»
\end{solution}
